\newpage
\section{Appendix}

\subsection{Equivalent results in $\mathbb{R}$}\label{sec:equiv-R}

\begin{remark}
    In $\mathbb{R}$, TFAE 
    \begin{itemize}
        \item Completeness Axiom (CA)
        \item Bolzano-Weierstrass (BW)
        \item Heine-Borel (HB)
        \item Nested Interval Property (NIP)
        \item Cauchy Convergence Criterion (CCC)
        \item Monotone Convergence Theorem (MCT)
    \end{itemize}
    This presents 30 potential proofs.
\end{remark}

The table below links to each implication. The entry in row $X$, column $Y$ is the proof of $X \Rightarrow Y$.

\begin{center}
\renewcommand{\arraystretch}{1.4}
\begin{tabular}{c|cccccc}
    $\Rightarrow$ & \textbf{CA} & \textbf{BW} & \textbf{HB} & \textbf{NIP} & \textbf{CCC} & \textbf{MCT} \\
    \hline
    \textbf{CA} & --- & \hyperref[pf:CA-BW]{\ref*{pf:CA-BW}} & \hyperref[pf:CA-HB]{\ref*{pf:CA-HB}} & \hyperref[pf:CA-NIP]{\ref*{pf:CA-NIP}} & \hyperref[pf:CA-CCC]{\ref*{pf:CA-CCC}} & \hyperref[pf:CA-MCT]{\ref*{pf:CA-MCT}} \\
    \textbf{BW} & \hyperref[pf:BW-CA]{\ref*{pf:BW-CA}} & --- & \hyperref[pf:BW-HB]{\ref*{pf:BW-HB}} & \hyperref[pf:BW-NIP]{\ref*{pf:BW-NIP}} & \hyperref[pf:BW-CCC]{\ref*{pf:BW-CCC}} & \hyperref[pf:BW-MCT]{\ref*{pf:BW-MCT}} \\
    \textbf{HB} & \hyperref[pf:HB-CA]{\ref*{pf:HB-CA}} & \hyperref[pf:HB-BW]{\ref*{pf:HB-BW}} & --- & \hyperref[pf:HB-NIP]{\ref*{pf:HB-NIP}} & \hyperref[pf:HB-CCC]{\ref*{pf:HB-CCC}} & \hyperref[pf:HB-MCT]{\ref*{pf:HB-MCT}} \\
    \textbf{NIP} & \hyperref[pf:NIP-CA]{\ref*{pf:NIP-CA}} & \hyperref[pf:NIP-BW]{\ref*{pf:NIP-BW}} & \hyperref[pf:NIP-HB]{\ref*{pf:NIP-HB}} & --- & \hyperref[pf:NIP-CCC]{\ref*{pf:NIP-CCC}} & \hyperref[pf:NIP-MCT]{\ref*{pf:NIP-MCT}} \\
    \textbf{CCC} & \hyperref[pf:CCC-CA]{\ref*{pf:CCC-CA}} & \hyperref[pf:CCC-BW]{\ref*{pf:CCC-BW}} & \hyperref[pf:CCC-HB]{\ref*{pf:CCC-HB}} & \hyperref[pf:CCC-NIP]{\ref*{pf:CCC-NIP}} & --- & \hyperref[pf:CCC-MCT]{\ref*{pf:CCC-MCT}} \\
    \textbf{MCT} & \hyperref[pf:MCT-CA]{\ref*{pf:MCT-CA}} & \hyperref[pf:MCT-BW]{\ref*{pf:MCT-BW}} & \hyperref[pf:MCT-HB]{\ref*{pf:MCT-HB}} & \hyperref[pf:MCT-NIP]{\ref*{pf:MCT-NIP}} & \hyperref[pf:MCT-CCC]{\ref*{pf:MCT-CCC}} & --- \\
\end{tabular}
\end{center}

% \subsubsection{Completeness Axiom $\Rightarrow$ Bolzano-Weierstrass}\label{pf:CA-BW}

% \begin{proof}
%     \remind{Exercise}
% \end{proof}

% \subsubsection{Completeness Axiom $\Rightarrow$ Heine-Borel}\label{pf:CA-HB}

% \begin{proof}
%     \remind{Exercise}
% \end{proof}

% \subsubsection{Completeness Axiom $\Rightarrow$ Nested Interval Property}\label{pf:CA-NIP}

% \begin{proof}
%     \remind{Exercise}
% \end{proof}

% \subsubsection{Completeness Axiom $\Rightarrow$ Cauchy Convergence Criterion}\label{pf:CA-CCC}

% \begin{proof}
%     \remind{Exercise}
% \end{proof}

% \subsubsection{Completeness Axiom $\Rightarrow$ Monotone Convergence Theorem}\label{pf:CA-MCT}

% \begin{proof}
%     \remind{Exercise}
% \end{proof}

% \subsubsection{Bolzano-Weierstrass $\Rightarrow$ Completeness Axiom}\label{pf:BW-CA}

% \begin{proof}
%     \remind{Exercise}
% \end{proof}

% \subsubsection{Bolzano-Weierstrass $\Rightarrow$ Heine-Borel}\label{pf:BW-HB}

% \begin{proof}
%     \remind{Exercise}
% \end{proof}

% \subsubsection{Bolzano-Weierstrass $\Rightarrow$ Nested Interval Property}\label{pf:BW-NIP}

% \begin{proof}
%     \remind{Exercise}
% \end{proof}

% \subsubsection{Bolzano-Weierstrass $\Rightarrow$ Cauchy Convergence Criterion}\label{pf:BW-CCC}

% \begin{proof}
%     \remind{Exercise}
% \end{proof}

% \subsubsection{Bolzano-Weierstrass $\Rightarrow$ Monotone Convergence Theorem}\label{pf:BW-MCT}

% \begin{proof}
%     \remind{Exercise}
% \end{proof}

% \subsubsection{Heine-Borel $\Rightarrow$ Completeness Axiom}\label{pf:HB-CA}

% \begin{proof}
%     \remind{Exercise}
% \end{proof}

% \subsubsection{Heine-Borel $\Rightarrow$ Bolzano-Weierstrass}\label{pf:HB-BW}

% \begin{proof}
%     \remind{Exercise}
% \end{proof}

% \subsubsection{Heine-Borel $\Rightarrow$ Nested Interval Property}\label{pf:HB-NIP}

% \begin{proof}
%     \remind{Exercise}
% \end{proof}

% \subsubsection{Heine-Borel $\Rightarrow$ Cauchy Convergence Criterion}\label{pf:HB-CCC}

% \begin{proof}
%     \remind{Exercise}
% \end{proof}

% \subsubsection{Heine-Borel $\Rightarrow$ Monotone Convergence Theorem}\label{pf:HB-MCT}

% \begin{proof}
%     \remind{Exercise}
% \end{proof}

% \subsubsection{Nested Interval Property $\Rightarrow$ Completeness Axiom}\label{pf:NIP-CA}

% \begin{proof}
%     \remind{Exercise}
% \end{proof}

% \subsubsection{Nested Interval Property $\Rightarrow$ Bolzano-Weierstrass}\label{pf:NIP-BW}

% \begin{proof}
%     \remind{Exercise}
% \end{proof}

% \subsubsection{Nested Interval Property $\Rightarrow$ Heine-Borel}\label{pf:NIP-HB}

% \begin{proof}
%     \remind{Exercise}
% \end{proof}

% \subsubsection{Nested Interval Property $\Rightarrow$ Cauchy Convergence Criterion}\label{pf:NIP-CCC}

% \begin{proof}
%     \remind{Exercise}
% \end{proof}

% \subsubsection{Nested Interval Property $\Rightarrow$ Monotone Convergence Theorem}\label{pf:NIP-MCT}

% \begin{proof}
%     \remind{Exercise}
% \end{proof}

% \subsubsection{Cauchy Convergence Criterion $\Rightarrow$ Completeness Axiom}\label{pf:CCC-CA}

% \begin{proof}
%     \remind{Exercise}
% \end{proof}

% \subsubsection{Cauchy Convergence Criterion $\Rightarrow$ Bolzano-Weierstrass}\label{pf:CCC-BW}

% \begin{proof}
%     \remind{Exercise}
% \end{proof}

% \subsubsection{Cauchy Convergence Criterion $\Rightarrow$ Heine-Borel}\label{pf:CCC-HB}

% \begin{proof}
%     \remind{Exercise}
% \end{proof}

% \subsubsection{Cauchy Convergence Criterion $\Rightarrow$ Nested Interval Property}\label{pf:CCC-NIP}

% \begin{proof}
%     \remind{Exercise}
% \end{proof}

% \subsubsection{Cauchy Convergence Criterion $\Rightarrow$ Monotone Convergence Theorem}\label{pf:CCC-MCT}

% \begin{proof}
%     \remind{Exercise}
% \end{proof}

% \subsubsection{Monotone Convergence Theorem $\Rightarrow$ Completeness Axiom}\label{pf:MCT-CA}

% \begin{proof}
%     \remind{Exercise}
% \end{proof}

% \subsubsection{Monotone Convergence Theorem $\Rightarrow$ Bolzano-Weierstrass}\label{pf:MCT-BW}

% \begin{proof}
%     \remind{Exercise}
% \end{proof}

% \subsubsection{Monotone Convergence Theorem $\Rightarrow$ Heine-Borel}\label{pf:MCT-HB}

% \begin{proof}
%     \remind{Exercise}
% \end{proof}

% \subsubsection{Monotone Convergence Theorem $\Rightarrow$ Nested Interval Property}\label{pf:MCT-NIP}

% \begin{proof}
%     \remind{Exercise}
% \end{proof}

% \subsubsection{Monotone Convergence Theorem $\Rightarrow$ Cauchy Convergence Criterion}\label{pf:MCT-CCC}

% \begin{proof}
%     \remind{Exercise}
% \end{proof}

\subsection{Series encountered and their convergence properties}

% Node and arrow styles shared by the two series diagrams below
\tikzset{
    srstmt/.style={draw, rounded corners=2pt, fill=blue!8, align=center,
        text width=4.2cm, minimum height=1cm, inner sep=4pt},
    sryes/.style={srstmt, fill=green!12},
    srno/.style={srstmt, fill=red!10},
    srunk/.style={srstmt, fill=gray!15},
    srask/.style={srstmt, fill=yellow!20, rounded corners=8pt},
    srtool/.style={draw, rounded corners=2pt, fill=blue!4, align=left, inner sep=5pt},
    srimp/.style={double, double distance=1.6pt, -{Implies[length=2.6mm]}},
    sreqv/.style={double, double distance=1.6pt,
        {Implies[length=2.6mm]}-{Implies[length=2.6mm]}},
    srarr/.style={-{Stealth[length=2.4mm]}, thick},
    srfail/.style={dashed, red!70!black, -{Stealth[length=2.2mm]}, thick},
    srlbl/.style={font=\small, fill=white, inner sep=1.5pt}
}

\subsubsection{Flow chart: which result applies to $\sum a_n$?}

\begin{remark}[Flow chart]
    \hfill
    \begin{center}
    \begin{adjustbox}{max width=\textwidth}
    \begin{tikzpicture}[node distance=1.0cm and 1.8cm]
        \node[srstmt] (s) {Given $\sum a_n$};
        \node[srask, below=of s] (q1) {Does $a_n \to 0$?};
        \node[srno, right=of q1] (d1) {$\sum a_n$ \textbf{diverges}\\
            {\small (contrapositive of: converges $\Rightarrow a_n \to 0$)}};
        \node[srask, below=of q1] (q2) {Is $a_n \geq 0$ for all $n$?};
        \node[srunk, right=of q2] (abs) {Apply the nonnegative branch below
            to $\sum |a_n|$};
        \node[sryes, below=of abs] (abs1) {$\sum |a_n|$ converges:\\
            $\sum a_n$ converges (absolutely)};
        \node[srunk, below=of abs1] (abs2) {$\sum |a_n|$ diverges:\\
            inconclusive for $\sum a_n$\\
            {\small (it may converge conditionally)}};
        \node[srstmt, below=1.6cm of q2] (nn) {Nonnegative series:\\
            $S_n$ is increasing};
        \node[srtool, below=2.8cm of nn] (tools) {
            \begin{tabular}{@{}l l l@{}}
                \textbf{Tool} & \textbf{Converges if} & \textbf{Diverges if} \\
                \hline
                Partial sums & $S_n$ bounded (or has a bounded subsequence)
                    & $S_n$ unbounded ($S_n \to \infty$) \\[3pt]
                Comparison & $a_n \leq b_n$ eventually, $\sum b_n$ converges
                    & $a_n \geq b_n \geq 0$ eventually, $\sum b_n$ diverges \\[3pt]
                Ratio ($a_n \neq 0$) & $\limsup \frac{a_{n+1}}{a_n} < 1$
                    & $\liminf \frac{a_{n+1}}{a_n} > 1$ \\[3pt]
                Root & $\limsup \sqrt[n]{a_n} < 1$
                    & $\limsup \sqrt[n]{a_n} > 1$ \\[3pt]
                Benchmarks & $\sum a^{n-1}$ for $|a| < 1$; $\sum \frac{1}{n^p}$ for $p > 1$
                    & $\sum a^{n-1}$ for $|a| \geq 1$; $\sum \frac{1}{n^p}$ for $p \leq 1$ \\
            \end{tabular}};
        \node[font=\small, below=2pt of tools, align=center]
            {For the ratio and root tests, a limsup or liminf equal to $1$ is inconclusive.};

        \draw[srarr] (s) -- (q1);
        \draw[srarr] (q1) -- node[srlbl, above] {No} (d1);
        \draw[srarr] (q1) -- node[srlbl, left] {Yes} (q2);
        \draw[srarr] (q2) -- node[srlbl, above] {No} (abs);
        \draw[srarr] (q2) -- node[srlbl, left] {Yes} (nn);
        \draw[srarr] (abs) -- (abs1);
        \draw[srarr] (abs1) -- (abs2);
        \draw[srarr] (nn) -- (tools);
    \end{tikzpicture}
    \end{adjustbox}
    \end{center}
\end{remark}

\subsubsection{Implication map}

\begin{remark}[Implications for general series]
    Here $\{a_n\} \subset \mathbb{R}$ has no sign condition.
    A double arrow is an implication. A dashed red arrow marked $\times$ is a
    converse that is \emph{false}.
    \begin{center}
    \begin{adjustbox}{max width=\textwidth}
    \begin{tikzpicture}
        \node[srstmt] (conv) at (0,0) {$\sum a_n$ converges};
        \node[srstmt, text width=6.2cm] (cauchy) at (8.6,0) {
            $\{S_n\}$ converges\\
            $\iff \{S_n\}$ is Cauchy\\
            $\iff \forall \varepsilon > 0\ \exists N:\
            m \geq n > N \Rightarrow \bigl|\sum_{k=n}^{m} a_k\bigr| < \varepsilon$};
        \node[srstmt] (abs) at (0,3) {$\sum |a_n|$ converges\\ (absolute convergence)};
        \node[srstmt] (zero) at (-2.7,-3.2) {$a_n \to 0$};
        \node[srstmt] (regroup) at (2.7,-3.2) {every regrouping $\sum b_k$
            converges, with the same sum};
        \node[srstmt] (lin1) at (8,-3.2) {$\sum a_n$ and $\sum b_n$ converge};
        \node[srstmt, text width=5.4cm] (lin2) at (8,-5.4) {
            $\sum (\alpha a_n + \beta b_n) = \alpha \sum a_n + \beta \sum b_n$};
        \node[font=\small, align=center, below=2pt of lin2]
            {($\sum a_n b_n = (\sum a_n)(\sum b_n)$ is false in general)};

        \draw[sreqv] (conv) -- node[srlbl, above] {def.\ + CCC} (cauchy);

        \draw[srimp] ([xshift=-0.9cm]abs.south) -- ([xshift=-0.9cm]conv.north);
        \draw[srfail] ([xshift=0.9cm]conv.north) -- ([xshift=0.9cm]abs.south)
            node[srlbl, midway, right=1pt] {$\times$\ conditional convergence};

        \draw[srimp] ([xshift=-1.8cm]conv.south) -- ([xshift=0.9cm]zero.north);
        \draw[srfail] ([xshift=1.7cm]zero.north) -- ([xshift=-0.9cm]conv.south)
            node[srlbl, midway, right=1pt] {$\times$};
        \draw[srimp] ([xshift=1.8cm]conv.south) -- ([xshift=-0.9cm]regroup.north);
        \draw[srfail] ([xshift=-1.7cm]regroup.north) -- ([xshift=0.9cm]conv.south)
            node[srlbl, midway, left=1pt] {$\times$};
        \draw[srimp] (lin1) -- (lin2);
    \end{tikzpicture}
    \end{adjustbox}
    \end{center}
    The converse failures are: the harmonic series ($a_n \to 0$ but diverges)
    and $\sum (-1)^n$ (regroups to $\sum 0 = 0$, but diverges).
    Absolute convergence is proved from the Cauchy criterion and the triangle
    inequality.
\end{remark}

\begin{remark}[Implications for nonnegative series]
    Here $a_n \geq 0$, so $\{S_n\}$ is increasing and the monotone convergence
    theorem applies.
    \begin{center}
    \begin{adjustbox}{max width=\textwidth}
    \begin{tikzpicture}
        % centre column: the two outcomes
        \node[sryes] (bc) at (0,0) {$\sum a_n$ converges};
        \node[srno] (bd) at (0,-6.2) {$\sum a_n$ diverges};
        \draw[dashed] (bc.south) -- (bd.north)
            node[srlbl, midway] {negations of each other};

        % right column: equivalent conditions
        \node[srstmt] (r1) at (7.2,2.4) {$S_n$ is bounded};
        \node[srstmt] (r2) at (7.2,0.8) {$S_n$ has a bounded (or convergent)
            subsequence};
        \node[srstmt] (r3) at (7.2,-0.9) {every regrouping $\sum b_k$ converges\\
            {\small ($\sum a_n = A \iff \sum b_k = A$)}};
        \node[srstmt] (r4) at (7.2,-2.5) {$\sum |a_n|$ converges};
        \node[srstmt] (r5) at (7.2,-6.2) {$S_n \to +\infty$};
        \foreach \r in {r1,r2,r3,r4} \draw[sreqv] (bc.east) -- (\r.west);
        \draw[sreqv] (bd.east) -- (r5.west);

        % left column: sufficient conditions (tests)
        \node[srstmt] (c1) at (-7.2,2.0) {$a_n \leq b_n$ eventually and
            $\sum b_n$ converges\\ {\small (Comparison)}};
        \node[srstmt] (c2) at (-7.2,0) {$\limsup \frac{a_{n+1}}{a_n} < 1$\\
            {\small (Ratio)}};
        \node[srstmt] (c3) at (-7.2,-2.0) {$\limsup \sqrt[n]{a_n} < 1$\\
            {\small (Root)}};
        \foreach \c in {c1,c2,c3} \draw[srimp] (\c.east) -- (bc.west);

        \node[srstmt] (d1) at (-7.2,-4.4) {$a_n \geq b_n \geq 0$ eventually and
            $\sum b_n$ diverges\\ {\small (Comparison)}};
        \node[srstmt] (d2) at (-7.2,-6.2) {$\liminf \frac{a_{n+1}}{a_n} > 1$\\
            {\small (Ratio)}};
        \node[srstmt] (d3) at (-7.2,-8.0) {$\limsup \sqrt[n]{a_n} > 1$\\
            {\small (Root)}};
        \foreach \d in {d1,d2,d3} \draw[srimp] (\d.east) -- (bd.west);

        \node[srtool, text width=12cm] at (0,-9.9) {
            \textbf{Benchmarks.}
            Geometric: $\sum a^{n-1}$ converges (to $\frac{1}{1-a}$) iff $|a| < 1$.
            $p$-series: $\sum \frac{1}{n^p}$ converges iff $p > 1$
            (harmonic series is $p = 1$).
            For the ratio and root tests, a limsup or liminf equal to $1$ is
            inconclusive.};
    \end{tikzpicture}
    \end{adjustbox}
    \end{center}
\end{remark}
